\section{Two manifold models of the same double cover}
\label{mod:section}

We now construct the two closed manifolds that will be distinguished
later.  We first replace the independent panel labels of the full
reflection by a quotient of their label group.  The quotient will retain
the covering and odd-degree nonexistence properties, while recording
enough information to rule out nontrivial symmetries of the triangulation
that preserve its label relations.  We then fill the doubled chambers
with boundary identifications depending on a linear function of the
chamber label.  Every choice gives a manifold model of the same space,
and the choice determines a subgroup of the geometric symmetries that
acts by homeomorphisms on that model.

Retain the finite flag-no-square PL triangulation \(\mathcal L\) of \(S\) chosen
at the end of Section~\ref{sec:reflection}, with vertex set
\(\mathcal V\) of cardinality at least \(5000\).  Write \(S_s\) for its
closed dual panel at \(s\).  Thus
\[
 I(x)=\{s\in\mathcal V:x\in S_s\}\quad(x\in S),
 \qquad I(x)=\varnothing\quad(x\in P\setminus S),
\]
and each nonempty \(I(x)\) is the vertex set of a simplex of \(\mathcal L\).
In particular, it has at most four elements.  The full label group is
\(\mathsf A_0=\Ftwo^{\mathcal V}\), with basis \(e_s\), and its
reflection is \(Q^0\).  We use the compatible finite CW structures on
the chamber and panels fixed in that section.

\subsection{A rigid space of label relations}

For \(v=\sum_s v_s e_s\in\mathsf A_0\), write
\[
 \supp(v)=\{s:v_s=1\},\qquad \wt(v)=|\supp(v)|.
\]
Choose a tetrahedron with vertices \(m_1,m_2,m_3,m_4\).  For
\(0\leq i\leq7\), choose pairwise disjoint subsets
\[
 T_i\subset\mathcal V\setminus\{m_1,m_2,m_3,m_4\},
 \qquad |T_i|=16\cdot2^i-1.
\]
There are enough vertices: these subsets require
\(16(2^8-1)-8=4072\) vertices outside the tetrahedron.  Define
\begin{equation}\label{mod:relation-space}
 c_i=e_{m_{\lfloor i/2\rfloor+1}}+\sum_{s\in T_i}e_s,
 \qquad
 \mathsf C=\langle c_0,\ldots,c_7\rangle\leq\mathsf A_0.
\end{equation}
The two vectors associated with \(m_j\) have precisely that vertex in
common support.  All other overlaps between these eight supports are
empty.

\begin{lemma}[Rigid label relations]\label{mod:rigid-relations}
The subspace \(\mathsf C\) has dimension eight, and every nonzero
element of \(\mathsf C\) has weight at least eight.  If a simplicial
automorphism \(\sigma\) of \(\mathcal L\) has
\(\sigma_*(\mathsf C)=\mathsf C\), where
\(\sigma_*(e_s)=e_{\sigma(s)}\), then \(\sigma\) is the identity.
\end{lemma}

\begin{proof}
For \(J\subset\{0,\ldots,7\}\), let \(r(J)\) be the number of pairs
\(\{2j-2,2j-1\}\), \(1\leq j\leq4\), contained in \(J\).
Cancellation occurs only at the marked vertex of such a pair, so
\begin{equation}\label{mod:weight-formula}
 \wt\left(\sum_{i\in J}c_i\right)
   =16\sum_{i\in J}2^i-2r(J),\qquad 0\leq2r(J)\leq8.
\end{equation}
For distinct \(J,J'\), the first terms on the right differ by at least
\(16\), while their correction terms differ by at most \(8\).  Thus
all subset sums have different weights.  Every nonempty subset sum has
weight at least \(16-8=8\), so the \(c_i\) are independent.

A coordinate permutation preserves weight.  Since each element of
\(\mathsf C\) has a different weight from the other elements, any
coordinate permutation preserving \(\mathsf C\) fixes every \(c_i\).
It consequently fixes each singleton
\[
 \supp(c_{2j-2})\cap\supp(c_{2j-1})=\{m_j\}.
\]
In particular, \(\sigma\) fixes the chosen tetrahedron pointwise.
Every triangular face of a closed triangulated three-manifold belongs
to exactly two tetrahedra.  If one of them is fixed pointwise, an
automorphism fixing that face must preserve the other tetrahedron and
fix its fourth vertex.  This propagates pointwise fixation across
faces.  The tetrahedron adjacency graph is connected: a path between
interior points of two tetrahedra can be put in PL general position
with respect to the triangulation, avoiding the one-skeleton and
crossing triangular faces one at a time.  Propagation along that graph
shows that every vertex of \(\mathcal L\) is fixed.
\end{proof}

Pass to the quotient labels
\begin{equation}\label{mod:label-quotient}
 \mathsf A=\mathsf A_0/\mathsf C,\qquad a_s=e_s\bmod\mathsf C.
\end{equation}
For \(J\subset\mathcal V\), put
\[
 \mathsf A_0(J)=\langle e_s:s\in J\rangle,\qquad
 \mathsf A(J)=\langle a_s:s\in J\rangle.
\]
The labels \(a_s\) span \(\mathsf A\).  They are nonzero and are
distinct for distinct \(s\), because a relation of weight one or two
cannot lie in \(\mathsf C\).  If \(J\) is a simplex of \(\mathcal L\), then
\begin{equation}\label{mod:simplex-independence}
 \mathsf C\cap\mathsf A_0(J)=0,
 \qquad \dim_{\Ftwo}\mathsf A(J)=|J|,
\end{equation}
since \(|J|\leq4\) and every nonzero element of \(\mathsf C\) has
weight at least eight.  Finally,
\(\dim_{\Ftwo}\mathsf A=|\mathcal V|-8>0\).

\subsection{The quotient reflection and its cyclic covers}

Use the quotient labels to define
\begin{equation}\label{mod:reflection}
 \begin{gathered}
 Q=(\mathsf A\times P)/\sim,\\
 (b,x)\sim(b',x')
 \ \Longleftrightarrow\
 x=x'\ \text{ and }\ b-b'\in\mathsf A(I(x)).
 \end{gathered}
\end{equation}
Write its points as \([b,x]\).  The label quotient gives a map
\[
 q:Q^0\longrightarrow Q,\qquad [b_0,x]\longmapsto[b_0\bmod\mathsf C,x].
\]

\begin{lemma}\label{mod:quotient-cover}
The map \(q\) is a regular finite covering with deck group
\(\mathsf C\).  The space \(Q\) is a connected finite aspherical CW
complex.  The folding map
\begin{equation}\label{mod:folding}
 r:Q\longrightarrow P,\qquad r([b,x])=x,
\end{equation}
is split by the inclusion of any chamber.
\end{lemma}

\begin{proof}
The group \(\mathsf C\) acts on \(Q^0\) by translation of its
\(\mathsf A_0\)-labels.  The stabilizer of \([b_0,x]\) in
\(\mathsf A_0\) is \(\mathsf A_0(I(x))\).  Equation
\eqref{mod:simplex-independence} shows that its intersection with
\(\mathsf C\) is trivial.  Hence \(\mathsf C\) acts freely.

The orbit relation for this action, written on
\(\mathsf A_0\times P\), identifies two points precisely when their
chamber coordinates agree and their label difference lies in
\(\mathsf C+\mathsf A_0(I(x))\).  Its image under the label quotient
is exactly the relation in \eqref{mod:reflection}.  Thus
\(Q^0/\mathsf C=Q\), with the quotient topologies.  A free action of a
finite group on a Hausdorff space is a covering action: for a point,
choose disjoint neighborhoods of it and each of its nontrivial
translates, then intersect the finitely many resulting neighborhoods
to obtain one whose translates are pairwise disjoint.  The space
\(Q^0\) is Hausdorff because it is a finite CW complex, so this applies
and proves the covering assertion.

For each open cell of \(P\), the set \(I(x)\) is constant on that cell,
because every panel is a subcomplex.  The relation in
\eqref{mod:reflection} therefore identifies entire copies of an open
cell by the identity.  The resulting finitely many open cells, with
the attaching maps induced from those of \(P\), give a finite CW
structure on \(Q\).  The chamber \(P\) is connected.  Each panel is
nonempty, and crossing \(S_s\) joins a chamber labeled \(b\) to one
labeled \(b+a_s\).  Since the \(a_s\) span \(\mathsf A\), all chambers
belong to one component.  Covering maps induce isomorphisms on all
higher homotopy groups.  The asphericity of \(Q^0\) from
Theorem~\ref{refl:input} thus gives the asphericity of \(Q\).
The formula for \(r\) is well-defined and continuous, and its
restriction to each chamber is the identity on \(P\).
\end{proof}

Set \(G=\pi_1(Q)\), choosing basepoints in a chamber.  Compose the
split epimorphism \(r_*:G\to H\) with the chamber character
\(w:H\to\mathbb Z\), and again denote the resulting epimorphism by
\(w:G\to\mathbb Z\).  For \(d\geq1\), let
\[
 G_d=w^{-1}(d\mathbb Z),\qquad Q_d\longrightarrow Q
\]
be the corresponding connected cyclic cover.  If
\(p_d:P_d\to P\) denotes the chamber cover, then \(Q_d\) is the
pullback of \(p_d\) along \(r\).  Indeed, its subgroup is
\(r_*^{-1}(H_d)=G_d\), and the pullback is connected because
\(G\to\mathbb Z/d\) is onto.  In particular, it has the chamber
description
\begin{equation}\label{mod:cyclic-chambers}
 \begin{gathered}
 Q_d=(\mathsf A\times P_d)/\sim_d,\\
 (b,y)\sim_d(b',y')
 \ \Longleftrightarrow\
 y=y'\ \text{ and }\
 b-b'\in\mathsf A(I(p_d(y))).
 \end{gathered}
\end{equation}
One way to verify this description, including its topology, is to map
\((b,y)\) to \(([b,p_d(y)],y)\) in \(Q\times_P P_d\).  Its fibers are
the displayed equivalence classes, and the induced continuous
bijection is a homeomorphism because its source is compact and the
pullback is Hausdorff.  Thus each chamber is \(P_d\); every lift of a
panel \(S_s\) retains the label change \(a_s\).

\begin{proposition}\label{mod:odd-models}
For every odd integer \(d\geq3\), the space \(Q_d\) has no closed
topological four-manifold model.
\end{proposition}

\begin{proof}
The fold for \(Q^0\) also has a chamber section.  Its connected cyclic
cover is consequently \(Q^0_d=Q^0\times_P P_d\).  Commutativity of
the folding maps gives
\[
 Q^0_d
   \cong Q^0\times_Q(Q\times_P P_d)
   =Q^0\times_Q Q_d.
\]
It follows that \(q\) pulls back to a finite covering
\(q_d:Q^0_d\to Q_d\); its total space is connected by the preceding
description of \(Q^0_d\).

Suppose that \(f:Z\to Q_d\) is a homotopy equivalence from a closed
connected topological four-manifold.  Under the isomorphism \(f_*\),
the subgroup defining \(Q^0_d\) determines a connected finite cover
\(Z'\to Z\).  Choose compatible basepoints and lift \(f\) to
\(f':Z'\to Q^0_d\).  A homotopy inverse to \(f\) and the two inverse
homotopies lift to the corresponding covers, so \(f'\) is a homotopy
equivalence.  The finite cover \(Z'\) is again a closed connected
topological four-manifold.  This contradicts
Proposition~\ref{refl:odd-models} for \(Q^0_d\).
\end{proof}

\subsection{Filling doubled chambers}

The next construction uses only local independence and spanning of
the panel labels.  We formulate it for arbitrary labels with these
properties so that the same proof also retains the double-cover model
for a full reflection.

Let \(K\) be a finite PL triangulation of \(S\), with vertex set
\(\mathcal U\) and closed dual panels \(S^K_u\).  Use the
panel-compatible chamber CW structures of
Section~\ref{sec:reflection} for this choice of \(K\), and their
lifts to \((P_2,S_2)\).  For \(y\in S_2\),
write \(\bar y=p_2(y)\) and set
\[
 J_K(y)=\{u\in\mathcal U:\bar y\in S^K_u\};
\]
set \(J_K(y)=\varnothing\) for \(y\in P_2\setminus S_2\).
Let \(\mathsf F\) be a finite-dimensional \(\Ftwo\)-vector space with
labels \(\lambda_u\in\mathsf F\), \(u\in\mathcal U\), which span
\(\mathsf F\) and are linearly independent on every simplex of \(K\).
For \(J\subset\mathcal U\), put
\(\mathsf F(J)=\langle\lambda_u:u\in J\rangle\).  The corresponding
doubled reflection is
\begin{equation}\label{mod:general-double-reflection}
 \begin{gathered}
 R_2=(\mathsf F\times P_2)/\sim_\lambda,\\
 (b,y)\sim_\lambda(c,y')
 \ \Longleftrightarrow\
 y=y'\ \text{ and }\ b-c\in\mathsf F(J_K(y)).
 \end{gathered}
\end{equation}
Here \(R_2\) depends on \(K,\mathsf F\), and its labels.  For
\(K=\mathcal L\), \(\mathsf F=\mathsf A\), and \(\lambda_s=a_s\), it is \(Q_2\)
by \eqref{mod:cyclic-chambers}.

Let
\[
 j:(V,S_2)\longrightarrow(P_2,S_2)
\]
be the relative filling equivalence of
Proposition~\ref{refl:chamber-data}; we identify \(\partial V=S_2\),
so \(j\) is the identity on the boundary.  Let \(\theta\) be the
nontrivial deck transformation of \(P_2\to P\).  It preserves \(S_2\),
and \(\overline{\theta y}=\bar y\).  In particular it preserves the
sets of lifted panels used in \eqref{mod:general-double-reflection}.

Fix a linear map \(\epsilon:\mathsf F\to\Ftwo\).  Take one copy \(V_b\)
of \(V\) for each \(b\in\mathsf F\), and use the homotopy equivalence
\begin{equation}\label{mod:twisted-chamber-maps}
 j_b=\theta^{\epsilon(b)}\circ j:V_b\longrightarrow P_2.
\end{equation}
Its boundary restriction is the homeomorphism
\(h_b=\theta^{\epsilon(b)}:\partial V_b=S_2\to S_2\).  In this
notation a point \(v\) in \(V_b\) means the point corresponding to
\(v\in V\) under the fixed identification of the copies.

Define \(\mathcal M_\epsilon\) as the quotient of
\(\coprod_{b\in\mathsf F}V_b\) by the following relation.  Interior
points are identified only with themselves.  For boundary points
\(v\in\partial V_b\) and \(v'\in\partial V_c\), set
\begin{equation}\label{mod:boundary-relation}
 \begin{gathered}
 (b,v)\approx_\epsilon(c,v')\\
 \Longleftrightarrow\
 h_b(v)=h_c(v')=:y
 \quad\text{and}\quad
 b-c\in\mathsf F(J_K(y)).
 \end{gathered}
\end{equation}
Across a lifted \(u\)-panel, the copies labeled \(b\) and
\(b+\lambda_u\) are glued on their identified boundaries by
\(v\mapsto\theta^{\epsilon(\lambda_u)}v\); at a meeting of panels,
\eqref{mod:boundary-relation} includes the corresponding combinations
of label changes.
For a fixed coordinate \(y\), the allowed label differences form a
subgroup, so this is an equivalence relation.  The chamber maps give
a continuous comparison map
\begin{equation}\label{mod:comparison-map}
 \varphi_\epsilon:\mathcal M_\epsilon\longrightarrow R_2,\qquad
 [b,v]\longmapsto[b,j_b(v)].
\end{equation}
It is well-defined because equivalent boundary points have the same
target coordinate and an allowed label difference.

We record the local shape of the panel decomposition.  This also
specifies the charts needed when \(V\) has no chosen triangulation.

\begin{lemma}[Panel coordinates]\label{mod:panel-coordinates}
Let \(y\in S_2\), and enumerate
\(J_K(y)=\{u_1,\ldots,u_q\}\).  Then \(1\leq q\leq4\).  A neighborhood
of \((y,0)\) in \(S_2\times[0,\infty)\) has coordinates in a
relative-open neighborhood of the origin in
\(\mathbb R^{4-q}\times[0,\infty)^q\) such that the lifted panel
\(p_2^{-1}(S^K_{u_i})\times\{0\}\) is the face \(x_i=0\).
The neighborhood may be chosen to meet no lifted panel with label
outside \(J_K(y)\).
\end{lemma}

\begin{proof}
The dual panels cover \(S\), and panels meet precisely over simplices,
which gives \(1\leq q\leq4\).  In a geometric simplex, the closed dual
block of a vertex consists of the points at which that vertex's
barycentric coordinate is maximal.  Let \(\sigma\) be the simplex of
\(K\) whose relative interior contains \(\bar y\).  At \(\bar y\)
the \(q\) coordinates indexed by \(J_K(y)\) are tied for the maximum;
all other coordinates are strictly smaller.  After shrinking a
neighborhood, only these \(q\) coordinates can be maximal.

The PL local product about a point of the relative interior of
\(\sigma\) has the directions in \(\sigma\) and the transverse cone
on its link.  Because \(K\) triangulates a PL three-manifold, that cone
supplies PL Euclidean transverse directions.  In this product the
barycentric coordinates belonging to \(\sigma\) acquire a common
positive scale factor in a transverse direction.  Their comparisons
are therefore determined by the coordinates within \(\sigma\).
Take the deviations of the \(q\) relevant coordinates from their
mean, denoted \(d_1,\ldots,d_q\), so that \(\sum_i d_i=0\).
These are \(q-1\) independent local coordinates in \(\sigma\); the
other simplex directions and the transverse directions together give
\(4-q\) coordinates.  In this chart the \(u_i\)-panel is given by
\(d_i=\max_j d_j\).  Lifting a sufficiently small neighborhood to
\(S_2\) gives the same description there.

Adjoin the inward collar coordinate \(t\geq0\).  On the factor
\(\{(d_i):\sum_i d_i=0\}\times[0,\infty)\), define
\begin{equation}\label{mod:orthant-coordinates}
 x_i=t+\max_j d_j-d_i,\qquad 1\leq i\leq q.
\end{equation}
This is a homeomorphism onto \([0,\infty)^q\).  Its continuous inverse
is
\[
 t=\min_i x_i,\qquad
 d_i=\frac1q\sum_jx_j-x_i.
\]
Moreover, \(x_i=0\) holds exactly when \(t=0\) and
\(d_i=\max_jd_j\).  Thus it identifies the panels with the indicated
zero faces.  The finitely many panels not containing \(y\) are closed,
so the initial neighborhood can be chosen disjoint from all of them.
\end{proof}

\begin{proposition}[Manifold models from doubled chambers]
\label{mod:models}
Let \(K\) be a finite PL triangulation of \(S\).  Let \(\mathsf F\)
be a finite-dimensional \(\Ftwo\)-vector space with labels
\((\lambda_u)_{u\in\mathcal U}\) spanning \(\mathsf F\) and linearly
independent on every simplex of \(K\).  Using the fixed relative
filling \(j:(V,S_2)\to(P_2,S_2)\), for every linear map
\(\epsilon:\mathsf F\to\Ftwo\) the space \(\mathcal M_\epsilon\)
defined by \eqref{mod:boundary-relation} is a closed connected
topological four-manifold, and
\(\varphi_\epsilon:\mathcal M_\epsilon\to R_2\) in
\eqref{mod:comparison-map} is a homotopy equivalence.  If \(R_2\) is
aspherical, then \(\mathcal M_\epsilon\) is aspherical.
\end{proposition}

\begin{proof}
For \(a\in\mathsf F\), the set of boundary coordinates allowing the
label difference \(a\) is
\begin{equation}\label{mod:allowed-coordinates}
 \begin{split}
 \Omega_a
   &=\{y\in S_2:a\in\mathsf F(J_K(y))\}\\
   &=\bigcup_{\substack{J\subset\mathcal U\\
                         \sum_{u\in J}\lambda_u=a}}
       \ \bigcap_{u\in J}p_2^{-1}(S^K_u).
 \end{split}
\end{equation}
The empty intersection is understood to be \(S_2\).  The equality
holds because an element of \(\mathsf F(J_K(y))\) is the sum of the
labels of some subset of \(J_K(y)\).  Each \(\Omega_a\) is a finite
union of closed panel intersections.  The relation between two fixed
chamber boundaries in \eqref{mod:boundary-relation} is the equality
of their \(h_b\)-coordinates over \(\Omega_{b-c}\), hence is closed.
On a single chamber the relation is its full diagonal, because
\(h_b\) is injective.  The finitely many closed relations between
distinct chambers, together with these diagonals, give a closed
equivalence relation on the
compact Hausdorff space \(\coprod_b V_b\).  Its quotient
\(\mathcal M_\epsilon\) is therefore compact Hausdorff, and each
chamber embeds in it.  The same argument with \(P_2\) in place of
\(V\) shows that \(R_2\) is compact Hausdorff and that its chambers
embed.  The panel-compatible CW structures also give \(R_2\) a
finite CW structure by the cell argument in
Lemma~\ref{mod:quotient-cover}.

The manifold \(V\) is connected, since \(j\) is a homotopy equivalence
and \(P_2\) is connected.  Every lifted panel is nonempty.
Consequently, for each \(u\) the chambers labeled \(b\) and
\(b+\lambda_u\) meet.  Since the labels span \(\mathsf F\), these
intersections connect all the chambers, proving connectedness.

Interior points already have four-manifold charts.  Consider a
boundary point with coordinate \(y\), and write
\(J_K(y)=\{u_1,\ldots,u_q\}\).  The labels of the chambers containing
its equivalence class form
\[
 b+\mathsf F(J_K(y))
   =\left\{b+\sum_{i=1}^q\eta_i\lambda_{u_i}:
                    (\eta_1,\ldots,\eta_q)\in\Ftwo^q\right\}.
\]
Local independence makes this parametrization bijective, so there are
exactly \(2^q\) chambers at the point.  Choose collars of their
boundaries, transporting the boundary coordinate to \(S_2\) by the
maps \(h_b\).  Lemma~\ref{mod:panel-coordinates} supplies the same
chart in each collar.  Shrink it to a relative-open product box
\[
 (-\rho,\rho)^{4-q}\times[0,\rho)^q
\]
that meets no panel outside \(J_K(y)\).  These boxes together form an
open saturated subset of the disjoint union of chambers: every allowed
label change within them belongs to \(\mathsf F(J_K(y))\).

In the chamber indexed by \(\eta\in\Ftwo^q\), map this box to
\(\mathbb R^{4-q}\times\mathbb R^q\) by
\[
 (z,x_1,\ldots,x_q)\longmapsto
       (z,(-1)^{\eta_1}x_1,\ldots,(-1)^{\eta_q}x_q).
\]
Two such images agree exactly on the faces on which the relevant
coordinates are zero, and these are precisely the identifications in
\eqref{mod:boundary-relation}.  The quotient of the boxes is therefore
\((-\rho,\rho)^4\).  More explicitly, the displayed map is a continuous
bijection after taking the quotient, and its inverse is continuous by
the pasting lemma on the finitely many closed orthants of
\((-\rho,\rho)^4\).  This gives a chart without boundary.  Thus
\(\mathcal M_\epsilon\) is locally Euclidean of dimension four.
Compactness supplies a finite cover by such charts and interior
charts, so their countable bases also give a countable basis for
\(\mathcal M_\epsilon\).  We have proved that it is a closed
topological four-manifold.

It remains to prove that \(\varphi_\epsilon\) is a homotopy
equivalence.  Enumerate the finitely many labels, and in both
\(\mathcal M_\epsilon\) and \(R_2\) take the images of their chambers
one at a time.  These images and their finite unions are compact,
hence closed.  When the new label is \(b\) and the earlier labels are
\(c\), its attaching set in the boundary coordinate \(S_2\) is
\[
 K_b=\bigcup_{\text{earlier }c}\Omega_{b-c}.
\]
Formula \eqref{mod:allowed-coordinates} expresses \(K_b\) as a finite
union of lifted panel intersections.  It is therefore a subcomplex
of the lifted boundary CW structure.  In the new source chamber the
attaching set is \(h_b^{-1}(K_b)\); in the new target chamber it is
\(K_b\).  The comparison on them is exactly the homeomorphism \(h_b\).
The quotient descriptions and the chamber embeddings show that
adjoining the new chamber is the pushout along these attaching sets.
For example, the natural map from that pushout to the corresponding
union is a continuous bijection from a compact space to a Hausdorff
space and hence is a homeomorphism.

The inclusions of \(K_b\) into \(S_2\), and of \(h_b^{-1}(K_b)\)
into \(\partial V_b\), are cofibrations because they are, up to the
boundary homeomorphism, inclusions of CW subcomplexes.  The inclusions
of the boundaries into \(P_2\) and \(V_b\) are cofibrations by their
collars.  Composing these cofibrations gives cofibrations from the
attaching sets into the new chambers on both sides.  Hence each
ordinary pushout is a homotopy pushout: its comparison with the double
mapping cylinder is a homotopy equivalence.  Homotopy pushouts
preserve a map of spans that is a homotopy equivalence at all three
objects.  Here the map on the attaching set is \(h_b\), the map on
the new chamber is the homotopy equivalence \(j_b\), and the map on
the preceding union is a homotopy equivalence by induction.  The
first chamber starts the induction.  This proves that
\(\varphi_\epsilon\) is a homotopy equivalence.  If \(R_2\) is
aspherical, lifting this equivalence to universal covers shows that
the universal cover of \(\mathcal M_\epsilon\) is contractible.
\end{proof}

\begin{corollary}[The double cover of a full reflection]
\label{refl:even-model}
For every finite flag triangulation \(K\) of \(S\), let \(Q^0(K)\)
be its full reflection, using one independent basis label for each
vertex.  The connected double cyclic cover of \(Q^0(K)\) defined by
the folding character has a closed connected aspherical topological
four-manifold model.
\end{corollary}

\begin{proof}
First suppose that the triangulation is PL relative to the fixed PL
structure on \(S\).  Take \(\mathsf F=\Ftwo^{\mathcal U}\), \(\lambda_u=e_u\), and
\(\epsilon=0\) in Proposition~\ref{mod:models}.  These labels span
and are independent on every simplex.  The target \(R_2\) is the
pullback description of the double cyclic cover of \(Q^0(K)\) from
Section~\ref{sec:reflection}.  The proposition supplies its
four-manifold model.  The full reflection is aspherical by
Theorem~\ref{refl:input}, so the cover and the model are aspherical.

For an arbitrary finite triangulation, its abstract complex is a
combinatorial three-manifold
\cite[Section~2.1, p.~3]{PaixaoSpreer2015}, and hence induces a PL
structure on \(S\).  By \cite[Theorem~2(2), p.~64]{Hamilton1976}, an
isotopy \(h_t:S\to S\) from the identity ends at a PL homeomorphism
from this structure to the fixed one.  The triangulation \(h_1(K)\)
has the same abstract flag complex and is PL for the fixed structure.
For this comparison, use the panel-compatible collared CW representative
of \(P\) for \(h_1(K)\) on both sides.  Extend \(h_t\) across its collar,
fixed away from that collar, to an isotopy with endpoint \(h_P\), and
use the pulled-back CW structure on the \(K\) side.  The relative
comparison maps from Section~\ref{sec:reflection} fix \(S\).  Since
\(\pi_1(S)\to H\) is surjective, they preserve the marking of \(H\)
and the folding character, so the induced reflected homotopy equivalences
lift to the corresponding cyclic covers.  Identify the full label spaces
by \(e_u\mapsto e_{h_1(u)}\).  With this identification, the map
\([a,x]\mapsto[a,h_P(x)]\) is a homeomorphism from \(Q^0(K)\) to
\(Q^0(h_1(K))\), since it carries each panel to its corresponding
panel.  It also identifies the cyclic covers: the folding maps commute
with \(h_P\), and an isotopy to the identity preserves the character
\(w\) on fundamental groups, up to an inner automorphism.  The PL case
therefore supplies a model for the double cover of \(Q^0(K)\) as well.
\end{proof}

\subsection{The two models and their actions}

Return to the fixed triangulation \(\mathcal L\) and the quotient labels
\((\mathsf A,a_s)\).  They satisfy the hypotheses of
Proposition~\ref{mod:models} by \eqref{mod:simplex-independence}.
For each linear \(\epsilon:\mathsf A\to\Ftwo\), denote the resulting
manifold \(\mathcal M_\epsilon\) by \(M_\epsilon\), and its comparison
map by
\[
 f_\epsilon:M_\epsilon\longrightarrow Q_2.
\]
The target is \(Q_2\) by \eqref{mod:cyclic-chambers}.  It is
aspherical as a cover of \(Q\), so every \(M_\epsilon\) is a closed
connected aspherical topological four-manifold.  Set
\begin{equation}\label{mod:common-group}
 \Pi=\pi_1(Q_2)=G_2,
\end{equation}
using compatible basepoints.  Each \(f_\epsilon\) gives a marking of
\(\pi_1(M_\epsilon)\) by \(\Pi\), well-defined up to inner
automorphism when basepoint paths are suppressed.

The group \(\mathsf A\times\Ftwo\) acts on \(Q_2\) by
\begin{equation}\label{mod:ambient-action}
 (a,t)[b,y]=[b+a,\theta^t y].
\end{equation}
The formula respects the relation in \eqref{mod:cyclic-chambers}
because \(\overline{\theta^t y}=\bar y\).  For a linear
\(\epsilon:\mathsf A\to\Ftwo\), define the graph subgroup
\begin{equation}\label{mod:graph-subgroup}
 \mathsf B_\epsilon
    =\{(a,\epsilon(a)):a\in\mathsf A\}
    \leq\mathsf A\times\Ftwo.
\end{equation}

\begin{proposition}[Actions on the manifold models]\label{mod:actions}
The action \eqref{mod:ambient-action} is faithful.  For each linear
\(\epsilon:\mathsf A\to\Ftwo\), its restriction to
\(\mathsf B_\epsilon\) is realized by a faithful action by
homeomorphisms on \(M_\epsilon\), and \(f_\epsilon\) is equivariant
for these actions.  Under the marking by \(\Pi\), the outer
automorphism induced by each of these homeomorphisms is the one
induced by the corresponding action on \(Q_2\).
\end{proposition}

\begin{proof}
An element of \(\mathsf A\times\Ftwo\) acting trivially on \(Q_2\)
must fix a point in a chamber interior.  Interior coordinates are
not identified between labels, and the nontrivial deck transformation
of \(P_2\) fixes no point.  The element therefore has both coordinates
zero, proving faithfulness.

For \(a\in\mathsf A\), define a map on the copies of \(V\) by changing
the label from \(b\) to \(b+a\) and using the identity on \(V\).
Its boundary coordinates satisfy
\[
 h_{b+a}(v)=\theta^{\epsilon(a)}h_b(v).
\]
Thus two previously identified boundary coordinates are both changed
by \(\theta^{\epsilon(a)}\).  Their projections to \(S\) stay the
same, and their label difference stays the same.  The map respects
\eqref{mod:boundary-relation}, and descends to a homeomorphism
\[
 T^\epsilon_a:M_\epsilon\longrightarrow M_\epsilon,\qquad
 T^\epsilon_a[b,v]=[b+a,v].
\]
The maps compose according to addition in \(\mathsf A\).  Their
action is faithful because an interior point in a chamber labeled
\(b\) is sent to an interior point in the distinct chamber labeled
\(b+a\) when \(a\ne0\).

Finally, linearity of \(\epsilon\) gives the exact equality
\[
 \begin{split}
 f_\epsilon T^\epsilon_a[b,v]
    &=[b+a,\theta^{\epsilon(b+a)}j(v)]\\
    &=(a,\epsilon(a))[b,\theta^{\epsilon(b)}j(v)]
      =(a,\epsilon(a))f_\epsilon[b,v].
 \end{split}
\]
This is the required equivariance.  Passing to fundamental groups
and suppressing the basepoint paths gives the stated equality of
outer automorphisms.
\end{proof}

Choose a nonzero linear functional \(\ell:\mathsf A\to\Ftwo\), which
exists because \(\dim_{\Ftwo}\mathsf A>0\), and set
\begin{equation}\label{mod:two-models}
 M=M_0,\qquad N=M_\ell.
\end{equation}
Their equivalences to \(Q_2\) make them homotopy equivalent.  The
homeomorphisms just constructed realize
\(\mathsf B_0=\mathsf A\times\{0\}\) on \(M\) and
\(\mathsf B_\ell\) on \(N\).  We identify \(\mathsf B_0\) with
\(\mathsf A\) in the remaining group-theoretic argument.
